\documentclass[11pt,a4paper]{appletmanual}
\usepackage[english]{babel}
\manuallang{en}

\appletname{Consumer Optimum}
\appletsub{Two goods, monotone preferences and a linear budget constraint. The
2D diagram and the 3D surface are the same thing seen two ways: move something
in either and the other follows.}
\appletdir{consumer-optimum}

\begin{document}
\manualtitlepage{Applet manual · Introductory Microeconomics · Universidad de los Andes\\
Translated from the GeoGebra applet “Maximización de Utilidad, Curvas de Nivel y…”.\\
MIT licensed. It opens in a browser; there is nothing to install.}

\manualtoc
\thispagestyle{empty}
\clearpage
\setcounter{page}{1}

% ==========================================================================
\section{What this is}

The consumer's problem is taught in two dimensions: indifference curves and a
budget line. But utility is a function of two variables, and indifference
curves are its level curves. The familiar diagram is the view from directly
above a surface.

This applet shows both views at once and keeps them coupled. Drag the bundle in
the plane and it moves on the surface; drag it on the surface and it moves in
the plane. The orthographic X--Z projection turns ``find the best affordable
bundle'' into ``find the top of this hill''.

\figher[1.0]{panels}{The two panels. Left: the textbook diagram, with the
utility map underneath. Right: the same function as a surface
$z = u(x, y)$, with the budget line lifted onto it.}

\begin{amkey}
The budget curve lifted onto the surface is the object that joins the two
views. Its highest point \emph{is} the optimum. Seeing it as a hill with a
summit takes the mystery out of the tangency condition.
\end{amkey}

\subsection{Who it is for}

The consumer-choice topic, and any course where the level curves of a function
of two variables need to stop being an abstract drawing. It also carries a
challenge mode with six preference families, scoring and saved progress, which
works well as homework.

%% Sections shared by every manual, English edition.
%% Pulled in with \input{common-en}. The only thing that differs between
%% applets is \appletfolder, which is also the folder name and the path on the site.

\section{Opening it}

There is nothing to install. The whole application is one \code{index.html}
file: it downloads no libraries, calls no server, and stores nothing outside
your browser.

\subsection{The three ways}

\begin{description}
\item[Double-click the file.] The quickest route. Find
  \code{\appletfolder/index.html} and open it. The browser shows it straight off the
  disk, with an address beginning \code{file://}. It works offline from the
  first moment.
\item[From the web.] If someone has published it, the address ends in
  \code{/\appletfolder/}. Once the page has loaded you can disconnect: it needs the
  network for nothing else.
\item[Served from a folder.] For a class, putting the folder on any static
  server is enough. With Python installed:
  \begin{quote}\code{python3 -m http.server 8000}\end{quote}
  then \code{http://localhost:8000/\appletfolder/} in the browser.
\end{description}

\begin{amnote}
Saving the page with \key{Ctrl}\,+\,\key{S} (or \key{Cmd}\,+\,\key{S} on a
Mac) leaves a copy of your own that still works offline and that can travel on a
memory stick or go out by email.
\end{amnote}

\subsection{Language and theme}

Two pairs of buttons sit at the top right.

\begin{ctrltable}{Control & What it does}
ES / EN & Switches the whole interface between Spanish and English on the spot.
          Nothing is recomputed, so you can switch mid-explanation without
          losing the state you had. \\
Dark / Light & Switches the theme. Dark is the original design; light exists
          for projectors and for printing. \\
\end{ctrltable}

Both choices are remembered in the browser, so next time it opens the way you
left it.

\subsection{Which browser}

Any reasonably recent one: Chrome, Edge, Firefox or Safari, on a computer,
tablet or phone. On a narrow screen the panels stack one above the other
instead of sitting side by side. No account, no permissions, no extensions.


% ==========================================================================
\section{The screen}

\fig[1.0]{interface}{The whole screen.}

\begin{description}
\item[Mode.] \ui{Explore} or \ui{Challenges}, top right.
\item[Preferences.] The $u(x, y)$ box and six shortcuts.
\item[Budget constraint.] $p_x$, $p_y$, $m$.
\item[Graphical method.] Which of the three ways of hunting for the optimum is
  active.
\item[Layers.] Fourteen, described below.
\item[Domain.] How far the axes run, with a button to fit them to the budget.
\end{description}

Under the two panels sits a slider whose meaning changes with the chosen
method, and at the foot of each panel a row of readouts.

\subsection{The readouts}

\begin{ctrltable}{Readout & What it is}
$x$, $y$ & The coordinates of the bundle $P$. \\
$u(P)$ & Utility at $P$. \\
Spend & What is spent against what is available, $p_x x + p_y y$ against
  $m$. \\
MRS & The marginal rate of substitution at $P$, $u_x/u_y$, reported
  positive. \\
$p_x/p_y$ & The relative price. At an interior optimum the two agree. \\
Gap & How far $P$ is from being the optimum, with a label: \emph{tangency},
  \emph{near} or \emph{far}. \\
max u & The utility of the true optimum. The number to reach. \\
\end{ctrltable}

\begin{amnote}
The MRS is reported positive, as $u_x/u_y$, and compared directly against
$p_x/p_y$. The GeoGebra original displayed $-u_x/u_y$ under a label reading
$u_x/u_y$, which made the two numbers impossible to compare by eye.
\end{amnote}

% ==========================================================================
\section{The three graphical methods}

\begin{description}
\item[Pick the bundle.] $P$ moves along the budget line, and the slider below
  runs it from $x = 0$ to $x = m/p_x$. The task is to climb while the colour
  under $P$ keeps getting lighter.
\item[Raise the level curve.] $P$ stays put and the slider raises the level
  $u$. The task is to find the level at which the curve stops cutting the line
  twice and stops cutting it at all: there it grazes at a single point.
\item[Free point.] $P$ moves anywhere in the quadrant, inside or outside the
  feasible set. Useful for talking about unaffordable bundles and about ones
  that leave money unspent.
\end{description}

Both panels are draggable under all three methods. In 3D the pointer is
ray-marched into the height field, so the bundle lands where the cursor
visually touches the terrain.

\fig[1.0]{optimum}{With the \ui{Show optimum} layer and the tangent switched
on.}

% ==========================================================================
\section{The camera}

Four buttons above the 3D panel. Three are true \textbf{orthographic}
projections --- no perspective divide --- so parallel edges stay parallel and
equal world steps are equal screen steps.

\begin{ctrltable}{Button & Projects onto, and which axes end up on screen}
X--Z & The plane $y = 0$: $x$ right, $u$ up. \\
X--Y & The plane $z = 0$: $x$ right, $y$ up. \\
Y--Z & The plane $x = 0$: $y$ right, $u$ up. \\
3D & None: free perspective, with orbit and zoom. \\
\end{ctrltable}

\fig[0.62]{xz}{The X--Z projection. The lifted budget curve projects to a
single arch whose peak is the optimum.}

\ui{X--Z} is the one to reach for: ``find the best affordable bundle'' becomes
``find the top of this hill''. \ui{X--Y} reproduces the 2D diagram exactly,
which makes the correspondence between the two panels concrete. In both side
views an indifference curve, being a level set, projects to a horizontal line.

\fig[0.62]{xy}{The X--Y projection. It is the 2D diagram, drawn by the 3D
engine.}

Dragging or arrow-keying away from a canonical view releases the button but
keeps the projection type. Dragging the bundle works in every view: the app
uses only the coordinates the active projection can express, so a horizontal
drag in \ui{Y--Z} moves along $y$ rather than along the collapsed $x$ axis.

% ==========================================================================
\section{The layers}

\begin{ctrltable}{Layer & What it draws}
Utility map & The field $u$ under a colour ramp. \\
Background contours & Level curves at regular intervals. \\
Level curve & The curve at the level the slider is set to. \\
Curve through P & The indifference curve through the bundle. \\
Feasible set & Veils what cannot be afforded. \\
3D slice & The budget line lifted onto the surface. \\
Tangent at P & The tangent to the indifference curve at $P$. \\
Gradient & $\nabla u$ at $P$. \\
Partials & $u_x$ and $u_y$ separately. \\
Cutting plane & The vertical plane over the budget line. \\
Floor grid & The ground grid in 3D. \\
Affordable only & Clips the surface to the feasible set. \\
Show optimum & Reveals the true optimum. \\
Leave a trail & Keeps the level curves swept so far. \\
\end{ctrltable}

% ==========================================================================
\section{Special cases}

The six preference shortcuts cover the cases a course needs. Two are worth
looking at slowly, because they break the tangency condition.

\fig[1.0]{complements}{Perfect complements, \code{min(2x,y)}. The indifference
curve has an elbow and there is no tangent to match.}

\fig[1.0]{substitutes}{Perfect substitutes, \code{2x+3y}. With
$p_x/p_y = 2/3$ the budget line and the indifference curves are parallel: the
whole line is optimal.}

\begin{amnote}
At the opening prices ($p_x = 2$, $p_y = 3$) the \code{2x+3y} shortcut gives
exactly $\text{MRS} = p_x/p_y$. It is the degenerate case where the consumer is
indifferent between every affordable bundle, and it is worth showing before it
turns up unannounced in an exam.
\end{amnote}

% ==========================================================================
\section{Challenge mode}

\fig[1.0]{challenge}{Challenge mode.}

Six levels, each with its own preference family and with prices and income
generated afresh on every attempt:

\begin{ctrltable}{Level & Family and task}
1 & Cobb--Douglas. Place $P$ on the optimal bundle. \\
2 & Asymmetric Cobb--Douglas. Raise the level curve until it grazes. \\
3 & Perfect substitutes. Place $P$. \\
4 & Perfect complements. Place $P$. \\
5 & Quasilinear. Raise the level curve. \\
6 & CES. Place $P$. \\
\end{ctrltable}

While a challenge is live, prices, income and the utility function are locked.
Checking an answer reveals the optimum, scores the attempt, and explains
\emph{which} optimality condition applied: tangency, a kink, or a corner.
Progress is kept in the browser.

\begin{amnote}
Alternating the task between ``place the bundle'' and ``raise the curve'' is
deliberate: they are the two graphical methods on the syllabus, and a student
who only practises one does not recognise the other in an exam.
\end{amnote}

% ==========================================================================
\section{Classroom activities}

\begin{activity}{The two views are one}
\amset Opening values ($u = x^{0.5}y^{0.5}$, $p_x = 2$, $p_y = 3$, $m = 24$).
Method \ui{Pick the bundle}.
\amexp Drag $P$ in the 2D panel and watch it move in the 3D one. Then drag it
in the 3D panel and watch it move in the 2D one.
\par\smallskip
Set the camera to \ui{X--Y}: the right panel becomes the left one. It is the
shortest proof that an indifference curve is a level curve.
\end{activity}

\begin{activity}{Finding the summit}
\amset Camera on \ui{X--Z}, method \ui{Pick the bundle}.
\amexp The lifted budget curve projects to an arch. Move the slider and watch
$u(P)$ rise to a maximum and fall again.
\par\smallskip
The \ui{max u} readout says $4.899$, which is $\sqrt{24}$. The optimum is at
$x^\ast = 6$, $y^\ast = 4$. Check it against the Cobb--Douglas formula
$x^\ast = \alpha m/p_x$ and see that the number matches the top of the arch.
\end{activity}

\begin{activity}{Tangency, found two ways}
\amset Method \ui{Pick the bundle}; then \ui{Raise the level curve}.
\amexp With the first method, climb the line while the colour under $P$ keeps
lightening: the \ui{Gap} readout falls until it says \emph{tangency}.
\par\smallskip
With the second, raise $u$ until the curve stops cutting the line twice. The
level at which it grazes is the same $4.899$. Both syllabus methods give the
same number, and the class should see it in the same minute.
\end{activity}

\begin{activity}{MRS against relative price}
\amset Opening values, \ui{Tangent at P} layer on.
\amexp With $P$ away from the optimum, \ui{MRS} $= 1.33$ and
$p_x/p_y = 0.67$: the tangent is steeper than the budget line. Move $P$ until
the two numbers agree.
\par\smallskip
There is the first-order condition, read as two numbers chasing each other.
Switch on \ui{Gradient} and \ui{Partials} to see where the MRS comes from.
\end{activity}

\begin{activity}{An optimum with no tangency}
\amset \ui{Complementarios / Complements} shortcut, \code{min(2x,y)}.
\amexp The \ui{max u} readout gives $6.000$, at $x^\ast = 3$, $y^\ast = 6$. But
the MRS at the elbow is undefined: the readout gives $0.00$ and no tangent
matches anything.
\par\smallskip
Ask how the optimum is found, then. The answer is that you solve $2x = y$
together with the budget line, and the tangency condition plays no part. It is
the case that breaks the recipe.
\end{activity}

\begin{activity}{Every bundle ties}
\amset \ui{Sustitutos / Substitutes} shortcut, \code{2x+3y}, opening prices.
\amexp $\text{MRS} = 2/3$ and $p_x/p_y = 2/3$: equal. The budget line and the
indifference curves are parallel, so \emph{any} bundle on the line is optimal.
\par\smallskip
Move $p_x$ slightly either way and the optimum jumps to an axis. That is corner
demand, and this is the only setting where you can watch it at the tipping
point.
\end{activity}

\begin{activity}{What cannot be afforded}
\amset Method \ui{Free point}. \ui{Feasible set} layer on.
\amexp Drag $P$ outside the veil. The \ui{Spend} readout passes $24.0$ and the
applet does not stop you: the bundle exists, it just cannot be paid for.
\par\smallskip
Then drag it inwards, leaving money unspent. With monotone preferences that is
never optimal, and you can see why: there is always a bundle to the north-east
with more colour on it.
\end{activity}

\begin{activity}{Challenges as homework}
\amset \ui{Challenges} mode.
\amexp All six levels generate fresh prices and income on every attempt, so the
answer cannot be memorised. On checking, the applet explains which condition
applied in that case.
\par\smallskip
It is a one-session assignment that marks itself. Progress is kept in the
student's browser, so ask for a screenshot of the score.
\end{activity}

% ==========================================================================
\section{Expression syntax}

\begin{ctrltable}{Element & What is allowed}
Operators & \code{+ - * / \^{} ( )} \\
Functions & \code{sqrt ln log log10 exp abs min max sin cos tan} \\
Constants & \code{e}, \code{pi} \\
Implicit multiplication & \code{2xy} is \code{2*x*y};
  \code{x\^{}(1/3)y\^{}(2/3)} parses as expected \\
Associativity & \code{\^{}} binds to the right and tighter than unary minus:
  \code{-x\^{}2} is $-(x^2)$ \\
\end{ctrltable}

% ==========================================================================
\section{How it is built}

Everything is hand-rolled so the page stays self-contained.

\begin{ctrltable}{Piece & Approach}
Expression language & Tokeniser $\rightarrow$ Pratt parser $\rightarrow$
  closure tree. No \code{eval} or \code{new Function}, so a strict content
  security policy cannot break it. \\
Partial derivatives & Symbolic differentiation over the tree, with constant
  folding. Falls back to central differences when the expression contains
  \code{min}, \code{max} or \code{abs} --- exactly the kink cases. \\
Level curves & Marching squares over a $141\times141$ sampled field, with
  saddle disambiguation and NaN-cell skipping. \\
Optimum & Dense scan then golden-section refinement along the budget line, then
  classification into interior, kink or corner. Robust where a first-order
  condition finds nothing. \\
3D surface & Painter's algorithm on canvas 2D. A height field sorts exactly by
  centroid depth, so no depth buffer is needed and there is no context-loss
  failure mode. \\
3D picking & Ray-march the unprojected pointer ray against the height field,
  then bisect. Crossings count in both directions, since in a side projection
  the ray runs level with the terrain and can surface from underneath it. \\
\end{ctrltable}

% ==========================================================================
\section{Changes from the GeoGebra original}

\begin{description}
\item[Colour ramp.] The original tinted level curves with three piecewise-linear
  channel functions (\code{redfun}, \code{greenfun}, \code{bluefun}). The
  replacement runs blue and green at the bottom, yellow and orange through the
  middle, then bright red into crimson and deep brown at the top. Because the
  top end darkens, height is not recoverable from lightness alone up there ---
  the colour bar and the numeric readouts carry that instead, and every curve
  drawn on the map gets a white casing so it stays legible against the dark
  browns.
\item[Marks.] The bundle, the budget line, the tangent, the gradient, the
  partials and the revealed optimum are fixed bright hues with a dark casing,
  none of them reusing a hue the ramp passes through.
\item[Colour range.] \code{uMIN}/\code{uMAX} were read off two domain corners.
  For $\ln(x) + y$ that lower corner is $-\infty$ and the ramp collapses. The
  rewrite takes the 2nd percentile of the sampled field.
\item[Domain.] The $x_C$, $y_C$ sliders reached $-10$, which puts
  \code{sqrt(x)} and \code{ln(x)} outside their domain. The quadrant is now
  clamped to $x, y \ge 0$.
\item[Budget parametrisation.] $x_s$ ranged over $[x_C, x_G]$ rather than
  $[0, m/p_x]$, so the ``bundle'' could sit at negative $x$ with $y$ above the
  budget line. It is now parametrised by $t \in [0, m/p_x]$.
\item[Prices.] $p_X$ and $p_Y$ sliders started at 0, making
  $y_s = (m - p_X x)/p_Y$ and the domain $m/p_X$ divide by zero. Both are now
  bounded below by $0.2$.
\item[MRS labelling.] The original displayed $-f_x/f_y$ under a label reading
  \code{RMS :=}\allowbreak\ \code{(du/dx)/(du/dy)}. MRS is now reported positive as $f_x/f_y$
  and
  compared directly against $p_x/p_y$, with the tangency gap as its own
  readout.
\item[Dead objects.] \code{afin}, \code{text1} (referring to
  $a_0, a_1, a_2$, which do not exist in the file), \code{BooleFuncion},
  \code{BooleTransparente}, \code{BooleMostrarVar},
  \code{BooleIntersecciones}, \code{BooleTangent}, \code{BooleTanPlane},
  \code{s}, \code{Impl}, \code{Boole} and \code{DX} were leftovers from the
  applet this was copied from. None were carried over.
\end{description}

% ==========================================================================
\section{Accessibility}

Both canvases are keyboard-operable: arrow keys move the bundle or the level,
and \key{+} and \key{-} zoom the 3D view. Light and dark themes are designed
separately, and both respond to the system setting. Numbers are tabular. Motion
respects \code{prefers-reduced-motion}.

% ==========================================================================
\section{Troubleshooting}

\begin{ctrltable}{Symptom & What is happening}
The map is all one colour & The function has a $-\infty$ inside the drawn
  domain. Raise the domain minimum or change the expression. \\
I cannot drag in 3D & Check you are not in \ui{Challenges} mode with the
  controls locked. Otherwise dragging works in all four views. \\
The surface comes out in pieces & The function is undefined over part of the
  quadrant. That is correct: there is no ground there. \\
The camera button switches itself off & You dragged away from the canonical
  view. The projection type is kept; press it again to recentre. \\
I lost my challenge score & It is kept in the browser. In a private window, or
  after clearing site data, it starts from zero. \\
\end{ctrltable}

%% Sections shared by every manual, English edition.
%% Pulled in with \input{common-en}. The only thing that differs between
%% applets is \appletfolder, which is also the folder name and the path on the site.

\section{Publishing it for a class}

Any static host will do: GitHub Pages, Netlify, a university web directory,
even a shared folder served over HTTP. The only thing to upload is the
\code{index.html} file, inside a folder named however you want the address to
read.

The repository ships a GitHub Actions workflow
(\code{.github/workflows/pages.yml}) that publishes every applet at once on each
push to the default branch. To switch it on, in the repository:
\emph{Settings\,$\rightarrow$\,Pages\,$\rightarrow$\,Source: GitHub Actions}.
That is the one step the workflow cannot do for itself, because creating a Pages
site needs administration rights that a \code{GITHUB\_TOKEN} is never granted.

\begin{amnote}
To hand the applet to students without reliable internet, send them the
\code{index.html} file directly. It weighs less than a photograph and opens with
a double-click.
\end{amnote}

\section{Licence and provenance}

The application code is MIT. Use it, change it and pass it on, in class or out
of it, with attribution.

This applet is a standalone rewrite of an original GeoGebra applet. The rewrite
is not a literal translation: where the original had a construction error, a
division by zero reachable with its own sliders, or a dead object inherited from
whatever applet it was copied from, this version fixes it and says so. The
section \emph{Changes from the GeoGebra original} lists those changes one by
one.

Everything is hand-written and dependency-free: the expression parser, the
symbolic differentiation, the level-curve tracing and the drawing. In
particular there is no \code{eval}, so what a student types into a text box
cannot become code, and a strict content security policy does not break the
page.


\end{document}
